%==============================================================================
\section{Categorical Data --- Dữ liệu phân loại trong học máy}
%==============================================================================

\subsection{Dữ liệu phân loại là gì?}

\begin{dinhnghia}[title={Categorical data}]
Dữ liệu phân loại trong ML là dữ liệu gồm các \textbf{nhóm} hoặc \textbf{nhãn}, thay vì giá trị số.
Các nhóm có thể:
\begin{itemize}
  \item \textbf{Nominal}: không có thứ tự tự nhiên (ví dụ màu sắc, giới tính).
  \item \textbf{Ordinal}: có thứ tự tự nhiên (ví dụ trình độ học vấn, nhóm thu nhập).
\end{itemize}
\end{dinhnghia}

\begin{ynghia}[title={Biểu diễn và one-hot}]
Dữ liệu phân loại thường biểu diễn bằng giá trị rời rạc (số nguyên hoặc chuỗi), và thường được mã hoá thành vector one-hot trước khi đưa vào mô hình.

One-hot encoding: với mỗi category tạo vector nhị phân có bit 1 ở vị trí tương ứng category đó và 0 ở các vị trí còn lại.
\end{ynghia}

\subsection{Kỹ thuật xử lý dữ liệu phân loại}

\begin{dongco}[title={Tiền xử lý bắt buộc với nhiều thuật toán}]
Xử lý dữ liệu phân loại là phần quan trọng của tiền xử lý ML vì nhiều thuật toán yêu cầu đầu vào số.
Tùy thuật toán và bản chất dữ liệu phân loại, có thể dùng: label encoding, ordinal encoding, binary encoding, \ldots
\end{dongco}

\begin{phuongphap}[title={Các kỹ thuật trong bài}]
One-Hot Encoding, Label Encoding, Frequency Encoding, Target Encoding, Binary Encoding.
\end{phuongphap}

\subsection{One-Hot Encoding}

\begin{dinhnghia}[title={One-hot encoding}]
Tạo vector nhị phân cho mỗi category; mỗi phần tử biểu diễn có/không có category đó.
Ví dụ biến màu với giá trị red, blue, green $\Rightarrow$ ba vector: $[1,0,0]$, $[0,1,0]$, $[0,0,1]$.
\end{dinhnghia}

\begin{vidu}[title={Ví dụ One-Hot Encoding với Pandas}]
\begin{lstlisting}
import pandas as pd

# Creating a sample dataset with a categorical variable
data = {'color': ['red', 'green', 'blue', 'red', 'green']}
df = pd.DataFrame(data)

# Performing one-hot encoding
one_hot_encoded = pd.get_dummies(df['color'], prefix='color')

# Combining the encoded data with the original data
df = pd.concat([df, one_hot_encoded], axis=1)

# Drop the original categorical variable
df = df.drop('color', axis=1)

# Print the encoded data
print(df)
\end{lstlisting}

\textbf{Output}:
\begin{lstlisting}
      color_blue    color_green    color_red
0        0              0              1
1        0              1              0
2        1              0              0
3        0              0              1
4        0              1              0
\end{lstlisting}
Tạo ba biến nhị phân \texttt{color\_blue}, \texttt{color\_green}, \texttt{color\_red} (1 nếu có màu tương ứng, 0 nếu không), dùng cho classification/regression.
\end{vidu}

\begin{luuy}[title={Hạn chế}]
One-hot phù hợp biến phân loại nhỏ, hữu hạn; với biến có cardinality lớn có thể tạo quá nhiều cột đặc trưng.
\end{luuy}

\subsection{Label Encoding}

\begin{dinhnghia}[title={Label encoding}]
Gán một giá trị số duy nhất cho mỗi category; thứ tự số có thể theo thứ tự category.
Ví dụ biến ``Size'' với small, medium, large $\Rightarrow$ có thể gán 0, 1, 2.
\end{dinhnghia}

\begin{vidu}[title={Ví dụ Label Encoding với scikit-learn}]
\begin{lstlisting}
from sklearn.preprocessing import LabelEncoder

# create a sample dataset with a categorical variable
data = ['small', 'medium', 'large', 'small', 'large']

# create a label encoder object
label_encoder = LabelEncoder()

# fit and transform the data using the label encoder
encoded_data = label_encoder.fit_transform(data)

# print the encoded data
print(encoded_data)
\end{lstlisting}

\textbf{Output}:
\begin{lstlisting}
[2 1 0 2 0]
\end{lstlisting}
Mặc định thứ tự mã hoá theo thứ tự alphabet; có thể đổi thứ tự bằng cách truyền danh sách tùy chỉnh vào \texttt{LabelEncoder}.
\end{vidu}

\begin{luuy}[title={Khi nào dùng label encoding}]
Hữu ích với biến ordinal (có thứ tự tự nhiên).
Với biến nominal cần thận trọng vì số có thể gợi ý thứ tự không tồn tại --- khi đó one-hot an toàn hơn.
\end{luuy}

\subsection{Frequency Encoding}

\begin{dinhnghia}[title={Frequency encoding}]
Thay mỗi category bằng \textbf{tần suất} (hoặc tỉ lệ) xuất hiện trong dataset.
Ý tưởng: category xuất hiện nhiều có thể quan trọng/hữu ích hơn cho thuật toán.
\end{dinhnghia}

\begin{vidu}[title={Ví dụ Frequency Encoding}]
\begin{lstlisting}
import pandas as pd

# create a sample dataset with a categorical variable
data = {'color': ['red', 'green', 'blue', 'red', 'green']}
df = pd.DataFrame(data)

# calculate the frequency of each category in the categorical variable
freq = df['color'].value_counts(normalize=True)

# replace each category with its frequency
df['color_freq'] = df['color'].map(freq)

# drop the original categorical variable
df = df.drop('color', axis=1)

# print the encoded data
print(df)
\end{lstlisting}

\textbf{Output}:
\begin{lstlisting}
      color_freq
0        0.4
1        0.4
2        0.2
3        0.4
4        0.4
\end{lstlisting}
Ví dụ: hai lần ``red'' và ba lần ``green'' trong tập 5 mẫu $\Rightarrow$ tần suất tương ứng 0.4 và 0.6.
\end{vidu}

\begin{luuy}[title={Ghi chú}]
Frequency encoding là lựa chọn thay one-hot/label khi cardinality cao; hiệu quả phụ thuộc dataset và thuật toán.
\end{luuy}

\subsection{Target Encoding}

\begin{dinhnghia}[title={Target encoding}]
Thay mỗi category bằng giá trị tổng hợp (thường là \textbf{trung bình}) của biến target (target) trong nhóm category đó.
Giúp nắm quan hệ giữa biến phân loại và target, có thể cải thiện dự đoán.
\end{dinhnghia}

\begin{vidu}[title={Ví dụ Target Encoding (label encoder + mean theo nhóm)}]
\begin{lstlisting}
import pandas as pd
from sklearn.preprocessing import LabelEncoder

# create a sample dataset with a categorical variable and a target variable
data = {'color': ['red', 'green', 'blue', 'red', 'green'],
   'target': [1, 0, 1, 0, 1]}
df = pd.DataFrame(data)

# create a label encoder object and fit it to the data
label_encoder = LabelEncoder()
label_encoder.fit(df['color'])

# transform the categorical variable using the label encoder
df['color_encoded'] = label_encoder.transform(df['color'])

# create a mean encoder object and fit it to the transformed data
mean_encoder = df.groupby('color_encoded')['target'].mean().to_dict()

# map the mean encoded values to the categorical variable
df['color_encoded'] = df['color_encoded'].map(mean_encoder)

# print the encoded data
print(df)
\end{lstlisting}

\textbf{Output}:
\begin{lstlisting}
   color     target     color_encoded
0  red        1           0.5
1  green      0           0.5
2  blue       1           1.0
3  red        0           0.5
4  green      1           0.5
\end{lstlisting}
\end{vidu}

\begin{luuy}[title={Lưu ý: data leakage}]
Target encoding rất mạnh nhưng dễ \textbf{rò rỉ nhãn} nếu tính thống kê trên toàn bộ tập (kể cả test). Trong thực hành: chỉ fit encoder trên tập train, hoặc dùng cross-validation/out-of-fold encoding.
\end{luuy}

\subsection{Binary Encoding}

\begin{dinhnghia}[title={Binary encoding}]
Mỗi category được gán mã nhị phân; mỗi chữ số cho biết category có mặt (1) hay không (0).
Mã nhị phân thường dựa trên vị trí category trong danh sách đã sắp xếp.
\end{dinhnghia}

\begin{vidu}[title={Ví dụ Binary Encoding với category\_encoders}]
\begin{lstlisting}
import pandas as pd
import category_encoders as ce

# create a sample dataset with a categorical variable
data = {'color': ['red', 'green', 'blue', 'red', 'green']}
df = pd.DataFrame(data)

# create a binary encoder object and fit it to the data
binary_encoder = ce.BinaryEncoder(cols=['color'])
binary_encoder.fit(df['color'])

# transform the categorical variable using the binary encoder
encoded_data = binary_encoder.transform(df['color'])

# merge the encoded variable with the original dataframe
df = pd.concat([df, encoded_data], axis=1)

# print the encoded data
print(df)
\end{lstlisting}

\textbf{Output}:
\begin{lstlisting}
   color    color_0    color_1
0   red       0           1
1   green     1           0
2   blue      1           1
3   red       0           1
4   green     1           0
\end{lstlisting}
\end{vidu}

\begin{luuy}[title={Phạm vi áp dụng}]
Binary encoding phù hợp biến có số category vừa phải; với số category rất lớn có thể kém hiệu quả.
\end{luuy}
